\section{2026 年 11 月核验增补}

本节按论文首次提交时间排列。每项方法附有可核验的外部入选依据。

% METHOD: 2605.07351
\paragraph{SplitGS-Loc：分离和重组高斯以消除二维至三维匹配歧义}
2026.05~\cite{2605_07351}。SplitGS-Loc 面向重建颜色时形成的高斯冗余和多对一像素匹配。它依据高斯属性组织用于定位的特征场，将含混匹配拆开，再以匹配后的二维至三维对应求相机位姿，减少单纯拟合光度特征给 PnP 带来的干扰。

论文在 Cambridge Landmarks 与室内定位数据评估平移和旋转误差，并区分有无 GS-CPR 精修。论文所说免逐场景训练指定位特征场，原有 RGB 高斯地图仍要训练。作者在论文链接的仓库提供训练、拆分和测试代码，本书未复现实验。

\textbf{入选依据} \href{https://github.com/misoleekr/splitgs-loc}{公开实现}已找到与论文对应的代码文件。

% METHOD: 2607.04127
\paragraph{LiDAR-GS-SLAM：在跟踪和建图之间复用局部协方差}
2026.07~\cite{2607_04127}。该方法将激光雷达扫描与用于渲染的二维高斯统一到球面域。跟踪时由广义 ICP 得到的局部协方差同时作为原语初始化、增密与裁剪的几何依据，建图修正后的协方差再给扫描到地图的配准提供约束。

作者在 KITTI Odometry、Oxford Spires 和 Newer College 上评价，长程车载序列和较密集手持采集分别用于考察不同条件。结果依赖所用扫描及几何匹配，没有测出所有激光雷达配置的实时界限。作者公开代码包含跟踪、建图与重建评估入口。

\textbf{入选依据} \href{https://github.com/lab-of-ai-and-robotics/lidar-gs-slam}{公开实现}已找到与论文对应的代码文件。

% METHOD: 2609.21347
\paragraph{Cube-Splat：将全景图像分面后共用一个位姿变量}
2026.09~\cite{2609_21347}。直接将全景图像送入针孔高斯渲染器会受到投影模型限制。Cube-Splat 将全景帧分解成固定朝向、共享光心的虚拟透视视图，并复用透视高斯渲染。每个分面产生的梯度用伴随变换映射到同一个位姿变量，因而不会把各分面误当成可以独立运动的相机。

作者在 SynPano、OmniBlender 和 PALVIO 共 31 个序列中报告轨迹误差和视图质量，并给出逐序列结果。该表示依赖准确的全景标定和固定分面朝向，拼接误差会影响接缝附近的光度一致性。系统针对大体静态的环境，未融合惯性或深度观测，缺纹理与重复结构下仍可能漂移。

\textbf{入选依据} 论文署名单位包括 Ant Group（\href{https://arxiv.org/abs/2609.21347}{论文原文}）。\href{https://github.com/guoxf304/cubesplat}{公开实现}已找到与论文对应的代码文件。

% METHOD: 2609.23491
\paragraph{Elevator-VIGS：把电梯轿厢运动与轿厢内相机运动分开}
2026.09~\cite{2609_23491}。相机在轿厢内看到的是相对于轿厢的运动，而惯导响应包含轿厢整体运动。Elevator-VIGS 在束调整中加入沿重力方向的运输状态，通过视觉语言和深度线索识别电梯事件，使轨迹与高斯地图能在移动坐标系下更新。

作者测试十六段真实电梯记录，并用模拟电梯序列补充过程真值。真实轨迹主要在上下电梯前后的控制点评价竖向误差，中间全程轨迹的证据来自仿真。设计针对沿重力方向运动的电梯。

\textbf{入选依据} 论文署名单位包括 Microsoft（\href{https://arxiv.org/abs/2609.23491}{论文原文}）。
